\chapter{2011年重庆卷（理）}

\section{单选题}

\begin{problem}
复数 \(\frac{\i^2 + \i^3 + \i^4}{1 - \i} =\)（\quad）

\choices
    {\(-\frac{1}{2} - \frac{1}{2}\i\)}
    {\(-\frac{1}{2} + \frac{1}{2}\i\)}
    {\(\frac{1}{2} - \frac{1}{2}\i\)}
    {\(\frac{1}{2} + \frac{1}{2}\i\)}
\begin{answer}
C
\end{answer}

\begin{solution}
由 \(\i^2+\i^3+\i^4=-1-\i+1=-\i\)，得
\(\displaystyle\frac{\i^2+\i^3+\i^4}{1-\i}=\frac{-\i}{1-\i}=\frac{(-\i)(1+\i)}{(1-\i)(1+\i)}=\frac{1-\i}{2}=\frac{1}{2}-\frac{1}{2}\i\)，故选 C.
\end{solution}
\end{problem}

\begin{problem}
“\(x < -1\)”是“ \( x^{2} - 1 > 0 \) ”的（\quad）

\choices
    {充分而不必要条件}
    {必要而不充分条件}
    {充要条件}
    {既不充分也不必要条件}
\begin{answer}
A
\end{answer}

\begin{solution}
因为 \(x<-1\)，所以 \(x^2>1\)，从而 \(x^2-1>0\)，故条件“\(x<-1\)”是条件“\(x^2-1>0\)”的充分条件. 但当 \(x=2\) 时，\(x^2-1>0\) 成立而 \(x<-1\) 不成立，所以它不是必要条件，故选 A.
\end{solution}
\end{problem}

\begin{problem}
已知 \(\lim\limits_{x\to +\infty}\left(\frac{2}{x - 1} + \frac{ax - 1}{3x}\right) = 2\)，则 \(a =\)（\quad）

\choices
    {\(-6\)}
    {2}
    {3}
    {6}
\begin{answer}
D
\end{answer}

\begin{solution}
因为 \(\displaystyle\lim_{x\to+\infty}\frac{2}{x-1}=0\)，且
\(\displaystyle\frac{ax-1}{3x}=\frac{a}{3}-\frac{1}{3x}\)，所以原极限等于 \(\frac{a}{3}\). 由 \(\frac{a}{3}=2\)，得 \(a=6\)，故选 D.
\end{solution}
\end{problem}

\begin{problem}
\((1 + 3x)^{n}\) （其中 \(n\in \N\) 且 \(n\geqslant 6\) ）的展开式中 \(x^{5}\) 与 \(x^{6}\) 的系数相等，则（\quad）

\choices
    {6}
    {7}
    {8}
    {9}
\begin{answer}
B
\end{answer}

\begin{solution}
展开式中 \(x^5\) 的系数为 \(\mathrm C_n^5 3^5\)，\(x^6\) 的系数为 \(\mathrm C_n^6 3^6\). 由题设，得
\(\mathrm C_n^5 3^5=\mathrm C_n^6 3^6\)，即
\(1=3\dfrac{\mathrm C_n^6}{\mathrm C_n^5}=3\cdot\dfrac{n-5}{6}\). 因而 \(n=7\)，故选 B.
\end{solution}
\end{problem}

\begin{problem}
下列区间中，函数 \( f(x) = |\ln (2 - x)| \) 在其上为增函数的是（\quad）

\choices
    {\((-\infty, 1]\)}
    {\(\left[-1, \frac{4}{3}\right]\)}
    {\(\left[0, \frac{3}{2}\right)\)}
    {\([1,2)\)}
\begin{answer}
D
\end{answer}

\begin{solution}
函数的定义域为 \(x<2\). 当 \(x\in(-\infty,1]\) 时，\(\ln(2-x)\geqslant0\)，所以 \(f(x)=\ln(2-x)\)，且 \(f'(x)=-\frac{1}{2-x}<0\)，函数在此区间上为减函数. 当 \(x\in[1,2)\) 时，\(\ln(2-x)\leqslant0\)，所以 \(f(x)=-\ln(2-x)\)，且 \(f'(x)=\frac{1}{2-x}>0\)，函数在此区间上为增函数. 选项 A 为减区间，B、C 均跨过 \(x=1\)，只有 D 是增区间，故选 D.
\end{solution}
\end{problem}

\begin{problem}
若 \(\triangle ABC\) 的内角 \(A\)，\(B\)，\(C\) 所对的边 \(a\)，\(b\)，\(c\) 满足 \((a + b)^2 - c^2 = 4\)，且 \(C = 60^\circ\)，则 \(ab\) 的值为（\quad）

\choices
    {\( \frac{4}{3} \)}
    {\(8 - 4\sqrt{3}\)}
    {1}
    {\( \frac{2}{3} \)}
\begin{answer}
A
\end{answer}

\begin{solution}
由余弦定理及 \(C=60^\circ\)，得 \(c^2=a^2+b^2-2ab\cos60^\circ=a^2+b^2-ab\). 因此
\((a+b)^2-c^2=a^2+2ab+b^2-(a^2+b^2-ab)=3ab=4\)，所以 \(ab=\frac{4}{3}\)，故选 A.
\end{solution}
\end{problem}

\begin{problem}
已知 \(a > 0\)，\(b > 0\)，\(a + b = 2\)，则 \(y = \frac{1}{a} + \frac{4}{b}\) 的最小值是（\quad）

\choices
    {\( \frac{7}{2} \)}
    {4}
    {\( \frac{9}{2} \)}
    {5}
\begin{answer}
C
\end{answer}

\begin{solution}
由柯西不等式，得
\(y=\dfrac{1^2}{a}+\dfrac{2^2}{b}\geqslant\dfrac{(1+2)^2}{a+b}=\frac{9}{2}\). 当且仅当 \(\dfrac{1}{a}=\dfrac{2}{b}\)，即 \(b=2a\) 时取等号. 联立 \(a+b=2\)，得 \(a=\frac{2}{3}\)，\(b=\frac{4}{3}\)，此时 \(y=\frac{9}{2}\)，故选 C.
\end{solution}
\end{problem}

\begin{problem}
在圆 \(x^{2} + y^{2} - 2x - 6y = 0\) 内，过点 \(E(0,1)\) 的最长弦和最短弦分别为 \(AC\) 和 \(BD\)，则四边形 \(ABCD\) 的面积为（\quad）

\choices
    {\( 5\sqrt{2} \)}
    {\(10\sqrt{2}\)}
    {\(15\sqrt{2}\)}
    {\(20\sqrt{2}\)}
\begin{answer}
B
\end{answer}

\begin{solution}
将圆的方程化为 \((x-1)^2+(y-3)^2=10\)，故圆心为 \(O(1,3)\)，半径为 \(\sqrt{10}\)，且 \(OE=\sqrt{(1-0)^2+(3-1)^2}=\sqrt{5}\). 过点 \(E\) 的最长弦是过圆心的直径，因此 \(AC=2\sqrt{10}\). 对于过 \(E\) 的弦，弦长为 \(2\sqrt{10-d^2}\)，其中 \(d\) 是圆心到弦的距离；当弦垂直于 \(OE\) 时，\(d\) 取最大值 \(OE=\sqrt{5}\)，所以最短弦 \(BD=2\sqrt{10-5}=2\sqrt{5}\). 此时 \(AC\perp BD\)，故
\(S_{ABCD}=\frac{1}{2}AC\cdot BD=\frac{1}{2}\cdot2\sqrt{10}\cdot2\sqrt{5}=10\sqrt{2}\)，故选 B.
\end{solution}
\end{problem}

\begin{problem}
高为 \(\frac{\sqrt{2}}{4}\) 的四棱锥 \(S - ABCD\) 的底面是边长为 1 的正方形，点 \(S\)，\(A\)，\(B\)，\(C\)，\(D\) 均在半径为 1 的同一球面上，则底面 \(ABCD\) 的中心与顶点 \(S\) 之间的距离为（\quad）

\choices
    {\( \frac{\sqrt{2}}{4} \)}
    {\( \frac{\sqrt{2}}{2} \)}
    {1}
    {\(\sqrt{2}\)}
\begin{answer}
C
\end{answer}

\begin{solution}
设正方形中心为 \(O\)，球心为 \(T\). 由正方形边长为 \(1\)，得 \(OA=\frac{\sqrt{2}}{2}\). 因为 \(A\)，\(B\)，\(C\)，\(D\) 在同一球面上，球心 \(T\) 在过 \(O\) 且垂直于底面的直线上，且 \(OT^2=1-OA^2=\frac{1}{2}\). 取底面为水平面，并令顶点 \(S\) 在其上方，若 \(T\) 在底面下方，则 \(ST\geqslant\frac{\sqrt{2}}{4}+\frac{\sqrt{2}}{2}>1\)，与 \(S\) 在半径为 \(1\) 的球面上矛盾，故 \(T\) 在底面上方，且 \(OT=\frac{\sqrt{2}}{2}\). 设 \(S\) 到底面的高为 \(h=\frac{\sqrt{2}}{4}\)，由 \(ST=1\) 得
\(OS^2=1+2h\cdot OT-OT^2=1+2\cdot\frac{\sqrt{2}}{4}\cdot\frac{\sqrt{2}}{2}-\frac{1}{2}=1\). 所以底面中心与顶点 \(S\) 的距离为 \(1\)，故选 C.
\end{solution}
\end{problem}

\begin{problem}
设 \(m\)，\(k\) 为整数，方程 \(mx^2 - kx + 2 = 0\) 在区间 \((0,1)\) 内有两个不同的根，则 \(m + k\) 的最小值为（\quad）

\choices
    {\(-8\)}
    {8}
    {12}
    {13}
\begin{answer}
D
\end{answer}

\begin{solution}
设方程的两个根为 \(r_1\)，\(r_2\)，其中 \(0<r_1<r_2<1\). 由韦达定理，\(r_1r_2=\frac{2}{m}>0\)，且 \(r_1+r_2=\frac{k}{m}>0\)，故 \(m>0\)，\(k>0\). 又因 \(r_1r_2<1\)，得 \(m>2\)，所以 \(m\geqslant3\). 方程有两个不同的实根，故判别式 \(\Delta=k^2-8m>0\). 另外，\(f(1)=m-k+2=m(1-r_1)(1-r_2)>0\)，所以 \(k<m+2\).

当 \(m=3\)，\(4\)，\(5\) 时，由 \(k<m+2\) 及 \(k\) 为正整数，分别有 \(k\leqslant4\)，\(5\)，\(6\)，但此时分别有 \(k^2\leqslant16<24=8m\)，\(k^2\leqslant25<32=8m\)，\(k^2\leqslant36<40=8m\)，均与 \(\Delta>0\) 矛盾. 当 \(m\geqslant7\) 时，\(k^2>8m\geqslant56\)，所以 \(k\geqslant8\)，从而 \(m+k\geqslant15\). 因此只需考察 \(m=6\)，此时 \(k^2>48\)，故 \(k\geqslant7\). 取 \(k=7\)，方程为 \(6x^2-7x+2=0=(3x-2)(2x-1)\)，两个根为 \(\frac{1}{2}\)，\(\frac{2}{3}\)，均在 \((0,1)\) 内. 所以 \(m+k\) 的最小值为 \(6+7=13\)，故选 D.
\end{solution}
\end{problem}

\section{填空题}

\begin{problem}
在等差数列 \(\{a_{n}\}\) 中，\(a_{3} + a_{7} = 37\)，则 \(a_{2} + a_{4} + a_{6} + a_{8} =\fillinblank{}\).
\begin{answer}
74
\end{answer}

\begin{solution}
等差数列中，等距离两项的和相等. 因为 \(a_3+a_7=2a_5=37\)，所以 \(a_2+a_8=2a_5=37\)，\(a_4+a_6=2a_5=37\). 因而 \(a_2+a_4+a_6+a_8=37+37=74\).
\end{solution}
\end{problem}

\begin{problem}
已知单位向量 \(e_1\)，\(e_2\) 的夹角为 \(60^\circ\)，则 \(|2e_1 - e_2| =\fillinblank{}\).
\begin{answer}
\(\sqrt{3}\)
\end{answer}

\begin{solution}
由 \(e_1\)，\(e_2\) 是单位向量且夹角为 \(60^\circ\)，得 \(e_1\cdot e_2=\cos60^\circ=\frac{1}{2}\). 因此
\(\lvert2e_1-e_2\rvert^2=4\lvert e_1\rvert^2+\lvert e_2\rvert^2-4e_1\cdot e_2=4+1-4\cdot\frac{1}{2}=3\)，从而 \(\lvert2e_1-e_2\rvert=\sqrt{3}\).
\end{solution}
\end{problem}

\begin{problem}
将一枚均匀的硬币抛掷 6 次，则正面出现的次数比反面出现的次数多的概率为\fillinblank{}.
\begin{answer}
\(\frac{11}{32}\)
\end{answer}

\begin{solution}
设正面出现次数为 \(X\)，则反面出现次数为 \(6-X\). 条件“正面出现的次数比反面出现的次数多”等价于 \(X>6-X\)，即 \(X=4\)，\(5\)，\(6\). 六次抛掷共有 \(2^6\) 种等可能结果，满足条件的结果数为 \(\mathrm C_6^4+\mathrm C_6^5+\mathrm C_6^6=15+6+1=22\)，所以所求概率为 \(\frac{22}{2^6}=\frac{11}{32}\).
\end{solution}
\end{problem}

\begin{problem}
已知 \(\sin \alpha = \frac{1}{2} + \cos \alpha\)，且 \(\alpha \in \left(0, \frac{\pi}{2}\right)\)，则 \(\frac{\cos 2\alpha}{\sin \left(\alpha - \frac{\pi}{4}\right)}\) 的值为\fillinblank{}.
\begin{answer}
\(-\frac{\sqrt{14}}{2}\)
\end{answer}

\begin{solution}
令 \(s=\sin\alpha\)，\(c=\cos\alpha\)，则 \(s-c=\frac{1}{2}\). 由 \(s^2+c^2=1\)，得 \(1-2sc=(s-c)^2=\frac{1}{4}\)，所以 \(sc=\frac{3}{8}\). 又因 \(\alpha\in\left(0,\frac{\pi}{2}\right)\)，有 \(s+c>0\)，从而 \((s+c)^2=1+2sc=\frac{7}{4}\)，即 \(s+c=\frac{\sqrt{7}}{2}\). 因此
\(\cos2\alpha=c^2-s^2=(c-s)(c+s)=-\frac{\sqrt{7}}{4}\)，且 \(\sin\left(\alpha-\frac{\pi}{4}\right)=\frac{s-c}{\sqrt{2}}=\frac{\sqrt{2}}{4}\). 所以原式等于 \(\frac{-\sqrt{7}/4}{\sqrt{2}/4}=-\frac{\sqrt{14}}{2}\).
\end{solution}
\end{problem}

\begin{problem}
设圆 \(C\) 位于抛物线 \(y^{2} = 2x\) 与直线 \(x = 3\) 所围成的封闭区域（包含边界）内，则圆 \(C\) 的半径能取到的最大值为\fillinblank{}.
\begin{answer}
\(\sqrt{6} - 1\)
\end{answer}

\begin{solution}
记封闭区域为 \(\Omega=\{(x,y):y^2\leqslant2x,\ x\leqslant3\}\). 该区域关于 \(x\) 轴对称且为凸集，因此若半径为 \(r\) 的圆面位于其中，关于 \(x\) 轴反射后再取两个圆心的中点，仍可得到位于区域内的同半径圆面. 故不妨设圆心为 \((h,0)\).

由直线 \(x=3\) 的限制，得 \(h+r\leqslant3\). 对圆周上的点写成 \(x=h+ru\)，\(y=rv\)，其中 \(u^2+v^2=1\). 要使圆面位于抛物线 \(y^2=2x\) 的右侧，需要对所有 \(-1\leqslant u\leqslant1\) 有
\(h+ru\geqslant\frac{r^2v^2}{2}=\frac{r^2(1-u^2)}{2}\)，即
\(h\geqslant\frac{r^2}{2}(1-u^2)-ru\). 当 \(r\geqslant1\) 时，右端关于 \(u\) 的最大值在 \(u=-\frac{1}{r}\) 处取得，最大值为 \(\frac{r^2+1}{2}\). 半径小于 \(1\) 时不可能达到下面构造的半径 \(\sqrt{6}-1>1\)，所以求最大值时可取 \(r\geqslant1\). 于是
\(\frac{r^2+1}{2}\leqslant h\leqslant3-r\)，从而 \(r^2+2r-5\leqslant0\)，故 \(r\leqslant\sqrt{6}-1\).

取 \(r=\sqrt{6}-1\)，\(h=3-r=4-\sqrt{6}\). 此时 \(\frac{r^2+1}{2}=4-\sqrt{6}=h\)，上述两个边界条件同时满足，所得圆面确实位于 \(\Omega\) 内. 因此半径能取到的最大值为 \(\sqrt{6}-1\).
\end{solution}
\end{problem}

\section{解答题}

\begin{problem}
设 \(a \in \R\)，\(f(x) = \cos x (a \sin x - \cos x) + \cos^2 \left(\frac{\pi}{2} - x\right)\) 满足 \(f\left(-\frac{\pi}{3}\right) = f(0)\)，求函数 \(f(x)\) 在 \(\left[\frac{\pi}{4}, \frac{11\pi}{24}\right]\) 上的最大值和最小值.
\end{problem}

\begin{solution}
\begin{enumerate}
\item 由 \(\cos^2\left(\frac{\pi}{2}-x\right)=\sin^2x\)，得 \(f(x)=a\sin x\cos x-\cos^2x+\sin^2x=\frac{a}{2}\sin 2x-\cos 2x\). 因此 \(f\left(-\frac{\pi}{3}\right)=-\frac{\sqrt{3}}{4}a+\frac{1}{2}\)，而 \(f(0)=-1\). 由题设条件得 \(-\frac{\sqrt{3}}{4}a+\frac{1}{2}=-1\)，解得 \(a=2\sqrt{3}\). 从而 \(f(x)=\sqrt{3}\sin 2x-\cos 2x=2\sin\left(2x-\frac{\pi}{6}\right)\). 令 \(t=2x-\frac{\pi}{6}\)，则当 \(x\in\left[\frac{\pi}{4},\frac{11\pi}{24}\right]\) 时，\(t\in\left[\frac{\pi}{3},\frac{3\pi}{4}\right]\). 在此区间内，\(\sin t\) 在 \(t=\frac{\pi}{2}\) 时取最大值 \(1\)，在 \(t=\frac{3\pi}{4}\) 时取最小值 \(\frac{\sqrt{2}}{2}\)，故 \(f(x)\) 的最大值为 \(2\)，最小值为 \(\sqrt{2}\).
\end{enumerate}
\end{solution}

\begin{problem}
某市公租房的房源位于 \(A\)，\(B\)，\(C\) 三个片区. 设每位申请人只申请其中一个片区的房源. 且申请其中任一个片区的房源是等可能的. 求该市的任 \(4\) 位申请人中：
\begin{enumerate}
\item 恰有 2 人申请 \(A\) 片区的房源的概率；
\item 申请的房源所在片区的个数 \(\xi\) 的分布列与期望.
\end{enumerate}
\end{problem}

\begin{solution}
\begin{enumerate}
\item 设事件“申请 \(A\) 片区的房源”为事件 \(E\)，则 \(P(E)=\frac{1}{3}\). 四位申请人的申请相互独立，所以恰有两人申请 \(A\) 片区的概率为 \(\mathrm C_4^2\left(\frac{1}{3}\right)^2\left(\frac{2}{3}\right)^2=\frac{8}{27}\).
\item \(\xi\) 的可能取值为 \(1\)，\(2\)，\(3\). 四人申请同一片区时 \(\xi=1\)，故 \(P(\xi=1)=3\left(\frac{1}{3}\right)^4=\frac{1}{27}\). 对于 \(\xi=2\)，先从三个片区中选出两个，再从这两个片区中进行申请且两个片区都出现，故 \(P(\xi=2)=\mathrm C_3^2\frac{2^4-2}{3^4}=\frac{14}{27}\). 对于 \(\xi=3\)，由容斥原理，四人申请三个片区且每个片区都出现的申请方式数为 \(3^4-\mathrm C_3^1 2^4+\mathrm C_3^2=36\)，所以 \(P(\xi=3)=\frac{36}{3^4}=\frac{4}{9}\). 分布列为

\begin{center}
\begin{tabular}{c|ccc}
\(\xi\) & \(1\) & \(2\) & \(3\) \\
\hline
\(P\) & \(\frac{1}{27}\) & \(\frac{14}{27}\) & \(\frac{4}{9}\)
\end{tabular}
\end{center}

因此 \(E(\xi)=1\cdot\frac{1}{27}+2\cdot\frac{14}{27}+3\cdot\frac{4}{9}=\frac{65}{27}\).
\end{enumerate}
\end{solution}

\begin{problem}
设 \( f(x) = x^3 + ax^2 + bx + 1 \) 的导数 \( f'(x) \) 满足 \( f'(1) = 2a \)，\( f'(2) = -b \)，其中常数 \(a\)，\(b \in \R\).
\begin{enumerate}
\item 求曲线 \( y = f(x) \) 在点 \((1, f(1))\) 处的切线方程；
\item 设 \( g(x) = f'(x) \e^{-x} \)，求函数 \( g(x) \) 的极值.
\end{enumerate}
\end{problem}

\begin{solution}
\begin{enumerate}
\item \(f'(x)=3x^2+2ax+b\). 由 \(f'(1)=2a\) 得 \(3+2a+b=2a\)，所以 \(b=-3\). 由 \(f'(2)=-b\) 得 \(12+4a+b=-b\)，代入 \(b=-3\) 解得 \(a=-\frac{3}{2}\). 因此 \(f(1)=1-\frac{3}{2}-3+1=-\frac{5}{2}\)，且 \(f'(1)=2a=-3\). 曲线在 \((1,f(1))\) 处的切线满足 \(y-f(1)=f'(1)(x-1)\)，即 \(y+\frac{5}{2}=-3(x-1)\)，整理得 \(6x+2y-1=0\).
\item 由上可得 \(f'(x)=3x^2-3x-3\)，所以 \(g(x)=(3x^2-3x-3)\e^{-x}\). 求导得 \(g'(x)=(-3x^2+9x)\e^{-x}=-3x(x-3)\e^{-x}\). 因为 \(\e^{-x}>0\)，所以当 \(x<0\) 时 \(g'(x)<0\)，当 \(0<x<3\) 时 \(g'(x)>0\)，当 \(x>3\) 时 \(g'(x)<0\). 因此 \(g(x)\) 在 \((-\infty,0)\) 上单调递减，在 \((0,3)\) 上单调递增，在 \((3,+\infty)\) 上单调递减，故 \(x=0\) 时取得极小值 \(g(0)=-3\)，\(x=3\) 时取得极大值 \(g(3)=\frac{15}{\e^3}\).
\end{enumerate}
\end{solution}

\begin{problem}
如图，在四面体 \(ABCD\) 中，平面 \(ABC \perp\) 平面 \(ACD\)，\(AB \perp BC\)，\(AD = CD\)，\(\angle CAD = 30^\circ\).
\begin{enumerate}
\item 若 \(AD = 2\)，\(AB = 2BC\)，求四面体 \(ABCD\) 的体积；
\item 若二面角 \(C - AB - D\) 为 \(60^{\circ}\)，求异面直线 \(AD\) 与 \(BC\) 所成角的余弦值.
\end{enumerate}

\bitmapfigure{img/2011/chongqing_science/q19_fig1.png}
\end{problem}

\begin{solution}
\begin{enumerate}
\item 设 \(F\) 为 \(AC\) 的中点. 因为 \(AD=CD\)，所以 \(DF\perp AC\). 又因为平面 \(ABC\perp\) 平面 \(ACD\)，且 \(DF\) 在平面 \(ACD\) 内，所以 \(DF\perp\) 平面 \(ABC\). 在直角三角形 \(AFD\) 中，\(AD=2\)，\(\angle CAD=30^\circ\)，故 \(DF=AD\sin30^\circ=1\)，\(AF=AD\cos30^\circ=\sqrt{3}\)，从而 \(AC=2\sqrt{3}\). 在直角三角形 \(ABC\) 中，设 \(BC=x\)，则 \(AB=2x\)，所以 \(AC^2=AB^2+BC^2=5x^2\)，解得 \(BC=\frac{2\sqrt{15}}{5}\)，\(AB=\frac{4\sqrt{15}}{5}\). 因此 \(S_{\triangle ABC}=\frac{1}{2}AB\cdot BC=\frac{12}{5}\)，四面体体积为 \(V=\frac{1}{3}S_{\triangle ABC}\cdot DF=\frac{4}{5}\).
\item 二面角和异面直线所成的角在相似变换下不变，故不妨令 \(AD=CD=2\). 取 \(F\) 为 \(AC\) 的中点，以 \(F\) 为原点建立空间直角坐标系，使 \(AC\) 在 \(y\) 轴上，平面 \(ABC\) 为 \(xOy\) 平面，则可设 \(A=(0,\sqrt{3},0)\)，\(C=(0,-\sqrt{3},0)\)，\(D=(0,0,1)\)，\(B=(x,y,0)\). 由 \(AB\perp BC\) 得 \((x,y-\sqrt{3},0)\cdot(-x,-\sqrt{3}-y,0)=0\)，即 \(x^2+y^2=3\). 平面 \(ABD\) 的一个法向量为 \(\bs{n}=\overrightarrow{AB}\times\overrightarrow{AD}=(y-\sqrt{3},-x,-\sqrt{3}x)\)，平面 \(ABC\) 的一个法向量为 \(\bs{k}=(0,0,1)\). 由二面角 \(C-AB-D=60^\circ\)，得 \(\frac{|\bs{n}\cdot\bs{k}|}{|\bs{n}|}=\cos60^\circ=\frac{1}{2}\)，即 \(12x^2=(y-\sqrt{3})^2+4x^2\)，所以 \((y-\sqrt{3})^2=8x^2\). 联立 \(x^2+y^2=3\)，得 \(9y^2-2\sqrt{3}y-21=0\)，所以 \(y=\sqrt{3}\) 或 \(y=-\frac{7\sqrt{3}}{9}\). 前者使 \(x=0\) 且 \(B=A\)，不合题意，故 \(y=-\frac{7\sqrt{3}}{9}\)，且 \(x^2=\frac{32}{27}\). 关于 \(y\) 轴反射不影响所求角，取 \(x=\frac{4\sqrt{6}}{9}\)，则 \(\overrightarrow{BC}=\left(-\frac{4\sqrt{6}}{9},-\frac{2\sqrt{3}}{9},0\right)\)，\(\overrightarrow{AD}=(0,-\sqrt{3},1)\)，从而 \(\overrightarrow{AD}\cdot\overrightarrow{BC}=\frac{2}{3}\)，\(|AD|=2\)，\(|BC|=\frac{2}{\sqrt{3}}\). 设异面直线 \(AD\) 与 \(BC\) 所成的锐角为 \(\theta\)，则 \(\cos\theta=\frac{|\overrightarrow{AD}\cdot\overrightarrow{BC}|}{|AD|\cdot|BC|}=\frac{\sqrt{3}}{6}\).
\end{enumerate}
\end{solution}

\begin{problem}
如图，椭圆的中心为原点 \(O\)，离心率 \(e = \frac{\sqrt{2}}{2}\)，一条准线的方程为 \(x = 2\sqrt{2}\).
\begin{enumerate}
\item 求该椭圆的标准方程；
\item 设动点 \(P\) 满足： \(\overrightarrow{OP} = \overrightarrow{OM} + 2\overrightarrow{ON}\)，其中 \(M\)，\(N\) 是椭圆上的点，直线 \(OM\) 与 \(ON\) 的斜率之积为 \(-\frac{1}{2}\). 问：是否存在两个定点 \(F_1\)，\(F_2\)，使得 \(|PF_1| + |PF_2|\) 为定值？若存在，求 \(F_1\)，\(F_2\) 的坐标；若不存在，说明理由.
\end{enumerate}

\bitmapfigure{img/2011/chongqing_science/q20_fig1.png}
\end{problem}

\begin{solution}
\begin{enumerate}
\item 设椭圆的长半轴为 \(a\)，半焦距为 \(c\)，则 \(e=\frac{c}{a}\)，右侧准线到原点的距离为 \(\frac{a}{e}\). 由题设 \(\frac{a}{e}=2\sqrt{2}\)，得 \(a=2\)，进而 \(c=ae=\sqrt{2}\)，所以 \(b^2=a^2-c^2=2\). 因此椭圆的标准方程为 \(\frac{x^2}{4}+\frac{y^2}{2}=1\).
\item 设 \(M=(x_1,y_1)\)，\(N=(x_2,y_2)\)，由 \(\overrightarrow{OP}=\overrightarrow{OM}+2\overrightarrow{ON}\)，可得 \(P=(x_1+2x_2,y_1+2y_2)\). 因为 \(M\)，\(N\) 在椭圆上，所以 \(x_1^2+2y_1^2=4\)，\(x_2^2+2y_2^2=4\). 又由直线 \(OM\) 与 \(ON\) 的斜率之积为 \(-\frac{1}{2}\)，得 \(\frac{y_1}{x_1}\cdot\frac{y_2}{x_2}=-\frac{1}{2}\)，即 \(x_1x_2+2y_1y_2=0\). 设 \(P=(x,y)\)，则
\[
x^2+2y^2=(x_1+2x_2)^2+2(y_1+2y_2)^2=\left(x_1^2+2y_1^2\right)+4\left(x_2^2+2y_2^2\right)+4\left(x_1x_2+2y_1y_2\right)=20
\]
所以 \(P\) 在椭圆 \(\frac{x^2}{20}+\frac{y^2}{10}=1\) 上. 该椭圆的长半轴为 \(2\sqrt{5}\)，半焦距为 \(\sqrt{20-10}=\sqrt{10}\)，其两个焦点为 \(F_1=(-\sqrt{10},0)\)，\(F_2=(\sqrt{10},0)\). 由椭圆的定义，\(|PF_1|+|PF_2|=2\cdot2\sqrt{5}=4\sqrt{5}\) 为定值，故存在这样的两个定点.
\end{enumerate}
\end{solution}

\begin{problem}
设实数数列 \(\{a_{n}\}\) 的前 \(n\) 项和 \(S_{n}\) 满足 \(S_{n + 1} = a_{n + 1}S_n\) （\(n\in \N^*\)）.
\begin{enumerate}
\item 若 \(a_{1}\)，\(S_{2}\)，\(-2a_{2}\) 成等比数列，求 \(S_{2}\) 和 \(a_{3}\)；
\item 求证：对 \(k \geqslant 3\) 有 \(0 \leqslant a_{k+1} \leqslant a_k \leqslant \frac{4}{3}\).
\end{enumerate}
\end{problem}

\begin{solution}
\begin{enumerate}
\item 由 \(S_1=a_1\) 及递推关系，得 \(S_2=a_2S_1=a_1a_2\). 若 \(S_2=0\)，则又因 \(S_2=a_1+a_2\)，可得 \(a_1=a_2=0\)，此时三项 \(a_1\)，\(S_2\)，\(-2a_2\) 均为 \(0\)，不能按等比数列定义构成等比数列，故 \(S_2\ne0\). 由三项成等比数列，得 \(S_2^2=-2a_1a_2=-2S_2\)，从而 \(S_2=-2\). 又由 \(S_3=S_2+a_3=a_3S_2\)，得 \(a_3=\frac{S_2}{S_2-1}=\frac{2}{3}\).
\item 对任意 \(n\geqslant2\)，由 \(S_n=S_{n-1}+a_n=a_nS_{n-1}\)，得 \((a_n-1)S_{n-1}=a_n\). 因为 \(a_n=1\) 会导致 \(0=1\)，所以 \(a_n\ne1\)，从而 \(S_{n-1}=\frac{a_n}{a_n-1}\)，并且 \(S_n=a_nS_{n-1}=\frac{a_n^2}{a_n-1}\). 再由 \(S_{n+1}=S_n+a_{n+1}=a_{n+1}S_n\) 可知 \(S_n\ne1\)，故 \(a_{n+1}=\frac{S_n}{S_n-1}=\frac{a_n^2}{a_n^2-a_n+1}\). 对 \(k\geqslant3\)，注意到 \(a_k^2-a_k+1=\left(a_k-\frac{1}{2}\right)^2+\frac{3}{4}>0\)，所以 \(a_{k+1}\geqslant0\). 从 \(a_3=\frac{2}{3}>0\) 出发，若 \(a_k\geqslant0\)，则
\[
a_k-a_{k+1}=a_k-\frac{a_k^2}{a_k^2-a_k+1}=\frac{a_k(a_k-1)^2}{a_k^2-a_k+1}\geqslant0
\]
由此归纳可得 \(0\leqslant a_{k+1}\leqslant a_k\leqslant a_3=\frac{2}{3}\leqslant\frac{4}{3}\)，即得所证结论.
\end{enumerate}
\end{solution}
